% ============================================================
%  CAPÍTULO 14 – Procedimientos de Seguimiento, Control y Calidad
% ============================================================
\chapter{Procedimientos de Seguimiento, Control y Calidad}

\section{Procedimiento de evaluación del proyecto}

La evaluación del progreso del proyecto se realizó en tres momentos clave:

\begin{enumerate}
  \item \textbf{Entrega del 1er trimestre (diciembre 2025)}: Se entregó la primera versión de la memoria con los apartados de análisis, diseño de la base de datos y los primeros módulos implementados (autenticación y resumen del día). El tutor revisó la documentación y proporcionó retroalimentación sobre la estructura y el alcance.
  \item \textbf{Entrega del 2do trimestre (febrero 2026)}: Se presentó una versión funcional intermedia de la aplicación con los módulos de clientes, rutinas y calendario operativos.
  \item \textbf{Entrega final (4 de mayo de 2026)}: Versión completa del proyecto con todos los módulos implementados, pruebas realizadas y memoria finalizada.
\end{enumerate}

\section{Indicadores de calidad}

Los siguientes indicadores se establecieron para medir la calidad de la solución:

\begin{itemize}
  \item \textbf{Cobertura funcional}: Porcentaje de requisitos funcionales implementados respecto a los definidos en el análisis. \textbf{Meta}: 100\,\%.
  \item \textbf{Tiempo de carga de pantallas}: Las pantallas con carga de datos remotos deben mostrar los resultados en menos de 2 segundos en condiciones normales de red.
  \item \textbf{Tasa de fallos en pruebas}: Número de casos de prueba fallidos / total. \textbf{Meta}: 0 fallos en pruebas al momento de la entrega.
  \item \textbf{Ausencia de credenciales expuestas}: Verificación de que ninguna clave de API ni credencial de Supabase está incluida en el repositorio público.
\end{itemize}

\section{Procedimiento de gestión de incidencias}

Durante el desarrollo, las incidencias se gestionaron de la siguiente forma:

\begin{enumerate}
  \item \textbf{Detección}: El desarrollador que detecta un error lo documenta en la descripción del commit o en un comentario en el código (\texttt{// TODO: fix}).
  \item \textbf{Registro}: Se crea una anotación en el chat del equipo o una \textit{issue} en GitHub con la descripción del problema, el módulo afectado y los pasos de reproducción.
  \item \textbf{Resolución}: El responsable del módulo afectado asume la corrección en la siguiente sesión de trabajo.
  \item \textbf{Validación}: Se vuelve a ejecutar el caso de prueba correspondiente para confirmar la corrección.
\end{enumerate}

\section{Procedimiento de gestión de cambios}

Los cambios en el alcance del proyecto (añadir o eliminar funcionalidades) debían acordarse entre ambos desarrolladores y comunicarse al tutor en la siguiente revisión. Los cambios menores (ajustes de UI, optimizaciones) no requerían aprobación formal.

\section{Participación de usuarios en la evaluación}

El prototipo de Figma fue mostrado a dos personas externas al equipo de desarrollo (potenciales usuarios del perfil de entrenador personal) antes del inicio de la implementación. Sus comentarios influyeron en la simplificación del flujo de creación de rutinas y en la priorización del módulo de mensajería.

\section{Cumplimiento del pliego de condiciones}

La siguiente tabla resume el grado de cumplimiento de los requisitos funcionales:

\begin{table}[H]
\centering
\begin{tabular}{@{}lcc@{}}
\toprule
\textbf{Requisito} & \textbf{Implementado} & \textbf{Notas} \\
\midrule
Autenticación (RF-01 a RF-03) & Sí & Completo \\
Resumen del Día (RF-04, RF-05) & Sí & Completo \\
Mis Clientes (RF-06 a RF-08) & Sí & Completo \\
Crear Rutina + API (RF-09 a RF-11) & Sí & Completo \\
Calendario (RF-12, RF-13) & Sí & Completo \\
Mensajes en tiempo real (RF-14 a RF-16) & Sí & Completo \\
\bottomrule
\end{tabular}
\caption{Grado de cumplimiento de requisitos funcionales.}
\end{table}

\section{Métricas de seguimiento por iteración}

Con el objetivo de disponer de un control cuantitativo del avance, se definieron métricas ligeras por iteración quincenal. Estas métricas no sustituyen a una metodología formal, pero aportan una visión objetiva del progreso y de la estabilidad del producto.

\begin{longtable}{@{}p{2.2cm}p{3.2cm}p{3cm}p{5.6cm}@{}}
\toprule
\textbf{Iteración} & \textbf{Métrica principal} & \textbf{Resultado} & \textbf{Interpretación} \\
\midrule
\endhead
I1 (Nov) & Historias completadas / planificadas & 6 / 8 & Avance inicial correcto; se replanifican dos historias de UI \\
I2 (Dic) & Incidencias abiertas al cierre & 5 & Carga asumible para periodo vacacional \\
I3 (Ene) & Tiempo medio de cierre de incidencia & 2.3 días & Buena capacidad de respuesta técnica \\
I4 (Feb) & Requisitos funcionales cerrados & 78\,\% & Se alcanza hito de versión intermedia funcional \\
I5 (Mar) & Casos de prueba en verde & 88\,\% & Madurez alta; pendientes ajustes de usabilidad \\
I6 (Abr) & Casos de prueba en verde & 100\,\% & Entrega final estable y validada \\
\bottomrule
\caption{Evolución de métricas de seguimiento por iteración.}
\end{longtable}

\section{Control de calidad documental}

Además del control del software, se aplicó control sobre la propia memoria técnica para asegurar trazabilidad y cumplimiento de la guía académica:

\begin{itemize}
  \item Revisión de coherencia entre requisitos (capítulo de análisis) y pruebas (capítulo de validación).
  \item Revisión cruzada de numeración de figuras, tablas y referencias internas.
  \item Comprobación de estilo uniforme (tipografía, márgenes, estructura de capítulos y pies de figura).
  \item Verificación de que cada decisión de diseño relevante tuviese justificación técnica explícita.
\end{itemize}

Este control permitió detectar desalineaciones tempranas, especialmente en la correspondencia entre capturas de pantalla y funcionalidades realmente implementadas.

\section{Lecciones de gestión del proyecto}

Desde la perspectiva de dirección técnica, se identificaron tres prácticas de alto impacto positivo:

\begin{enumerate}
  \item \textbf{Cerrar verticales funcionales}: completar un flujo extremo a extremo (UI + lógica + BD + prueba) antes de iniciar otro módulo.
  \item \textbf{Revisiones cortas y frecuentes}: validar cada avance en sesiones breves evita acumulación de deuda técnica.
  \item \textbf{Documentación incremental}: redactar la memoria en paralelo al desarrollo mejora precisión y reduce trabajo de última hora.
\end{enumerate}

Estas prácticas se consideran transferibles a futuros proyectos académicos y profesionales, especialmente en equipos pequeños como el del presente TFG.
